% Ejemplo de uso en Overleaf de las tablas y gráficos generados por model/metricas.py
% Subir a Overleaf la carpeta docs/metricas/ y agregar al preámbulo:
\documentclass{article}
\usepackage[utf8]{inputenc}
\usepackage[spanish]{babel}
\usepackage{booktabs}
\usepackage{pgfplots}
\pgfplotsset{compat=1.17}
\usepackage[margin=1.5cm,landscape]{geometry}
\begin{document}
\input{metricas/estrategias.tex}
\input{metricas/por_tick.tex}
\input{metricas/rondas.tex}
\input{metricas/amdahl.tex}
\input{metricas/ops_roofline.tex}
\input{metricas/recursos.tex}
\begin{figure}[htbp]
  \centering
  \begin{minipage}{0.48\linewidth}\input{metricas/roofline_donchian.tex}\end{minipage}\hfill
  \begin{minipage}{0.48\linewidth}\input{metricas/roofline_mul.tex}\end{minipage}
  \caption{Roofline del Espino (25\,MHz, CPI mínimo 2) con el kernel de cada versión.}
\end{figure}
\end{document}
